\chapter*{Abstract}
\addcontentsline{toc}{chapter}{Abstract}

This Master 1 computer science project presents the development of an automatic valuation
system for residential real estate (houses and apartments) in Senegal, based on machine
learning techniques. The goal is twofold: to compare several prediction models through
systematic hyperparameter optimization, and to make the results accessible through an
interactive graphical interface.

In the absence of any existing Senegalese real-estate dataset, a corpus of 660 sale listings
was collected through automated web scraping from the Expat-Dakar platform, covering the
whole country. After cleaning (removal of non-exploitable prices and of properties outside the
single-unit residential scope), a dataset of 603 listings was retained, described by surface
area, number of bedrooms, property type, and neighborhood.

Four models were trained and compared: a linear regression (baseline), a Random Forest, a
gradient boosting model (XGBoost), and a multi-layer perceptron (MLP) neural network — the
latter three tuned through randomized hyperparameter search with cross-validation. The
optimized Random Forest achieved the best performance on the test set (MAE of 73.3 million
FCFA, MAPE of 34.0\%, $R^2$ of 0.751), outperforming both XGBoost and the neural network — a
result consistent with the literature on modestly sized tabular datasets.

An interactive web application built with Streamlit allows users to estimate a property's
price from its characteristics, explore the dataset, and visually compare model performance.

\vspace{0.5cm}
\noindent\textbf{Keywords:} real estate, machine learning, hyperparameter optimization, Random
Forest, XGBoost, neural network, Senegal, Streamlit.
